\SourceLeaf{494}%
the theodolites made by the English instrument maker \textsc{Troughton} are the most convenient and useful. Prof. \textsc{Schumacher} has been kind enough to sell me his for this purpose at the price he paid, as he hopes to obtain another from England by the time he needs one himself again this year. — Some angle measurements will scarcely be practicable except at night, using Argand lamps with parabolic reflectors, which can be made visible at very great distances. To ensure that these are available when needed, I have ordered several such reflectors from a very skilled maker, the court instrument maker \textsc{Körner} in Jena, who has also supplied similar work for the Prussian surveys. He has promised to complete them by Easter. I leave unmentioned several other pieces of apparatus discussed in advance.

With all these arrangements made, I now believe I can begin work this coming spring. If the instrument makers duly keep all their promises, the work can then continue to advance rapidly. In any event, the various preliminary tasks, traveling through the regions, selecting the stations, building signal towers, making preliminary measurements and the like — (I already have the aforementioned theodolite by \textsc{Troughton}) — will first fill several months. Even in the unfavorable event that the larger \textsc{Reichenbach} theodolite we are expecting has still not been delivered by the time the preliminary work has advanced far enough for the measurements themselves to begin, the work need not come to a standstill. If necessary, I could begin with the angle measurements for which the observatory’s other equipment suffices. Indeed, in the very worst case, which I hope will not arise, of the theodolite being delayed still longer, I would make only a change in the order of work. I would use the rest of the summer for the astronomical work with the zenith sector here in Göttingen, which would otherwise be reserved for a later time if the entire summer can be devoted to triangulation.

Göttingen, March 20, 1821.
\begin{flushright}C. F. \textsc{Gauss}.\end{flushright}

\par\bigskip\Needspace{13\baselineskip}
\pdfbookmark[0]{Errata.}{druckfehler}
\begin{center}\large ERRATA.\end{center}
\par\noindent On page 296, for
\begin{center}
\begin{tabular}{c|c|c|c}
\(Q+q\)&\(P+p\)&\(\log m\)&\(k\)\\
&&\(+\)&\\
\(52^{0}\ 30'\)&\(52^{0}\ 32'\ 1''78428\)&&\(0''006\)
\end{tabular}
\end{center}
\noindent read
\begin{center}
\begin{tabular}{c|c|c|c}
\(Q+q\)&\(P+p\)&\(\log m\)&\(k\)\\
&&\(-\)&\\
\(52^{0}\ 30'\)&\(52^{0}\ 32'\ 1''78428\)&&\(0''006\)
\end{tabular}
\end{center}
